\documentclass[12pt,english]{article}
\usepackage[T1]{fontenc}
\usepackage[latin9]{inputenc}
\usepackage[toc,page]{appendix}
\usepackage{geometry}
\geometry{verbose,tmargin=1.5in,bmargin=1.5in,lmargin=1in,rmargin=1in}
\usepackage{bm}
\usepackage{amstext}
\usepackage{comment}
\usepackage{amsmath}
\usepackage{hyperref}
\hypersetup{
     colorlinks = true,
     citecolor = navy,
     linkcolor = navy,
     urlcolor=navy
}
\usepackage[space]{grffile}
\usepackage{setspace}
\usepackage[normalem]{ulem}
\usepackage{xcolor}
\usepackage{soul}
\usepackage{dsfont}
\usepackage{refstyle}
\usepackage[capitalize,nameinlink]{cleveref}
\usepackage{setspace}
\usepackage{bbm}
\usepackage[tableposition=t]{caption} % ensures caption of longtables is normalsized.
\usepackage{mathpazo} % add Palatino font
\usepackage{longtable}
\usepackage{amssymb}
\usepackage{amsfonts}
\usepackage{amsmath}
%\usepackage[margin=1in]{geometry} % to set margins
\usepackage{scalefnt}
\usepackage{wrapfig}
\usepackage{graphicx}
%\usepackage{subfigure}
\usepackage{lscape}
\usepackage{pdflscape}
\usepackage{multirow}
\usepackage{array}
\usepackage{rotating}
\usepackage{bibunits}
\usepackage{fix-cm}
\usepackage{float}
\floatplacement{table}{H}
\usepackage{epsfig}
\usepackage{epstopdf}
\usepackage{refstyle}
\usepackage{booktabs}
\usepackage{tabularx}
\usepackage{array}

\newcolumntype{x}[1]{>{\raggedleft\arraybackslash}p{#1}}%
\newcolumntype{d}[1]{>{\centering\arraybackslash}p{#1}}%
\newcolumntype{y}[1]{>{\raggedright\arraybackslash}p{#1}}%
\usepackage{verbatim}
%\epstopdfsetup{suffix=}
\setcounter{MaxMatrixCols}{10}
\usepackage[longnamesfirst,authoryear]{natbib}
\newenvironment{proof}[1][Proof]{\noindent \textbf{#1.}}{\rule{0.5em}{0.5em}}
\setlength{\topmargin}{-0.80in} \setlength{\textheight}{9.50in}
\setlength{\evensidemargin}{.0in} \setlength{\oddsidemargin}{.0in}
\setlength{\textwidth}{6.85in}
\setlength{\voffset}{0.4cm}
\setlength{\hoffset}{-0.5cm}
\usepackage{setspace}
\usepackage{afterpage}
\usepackage[flushleft]{threeparttable}
\usepackage{tabulary}
%\usepackage[titletoc,title]{appendix}
\usepackage[bottom]{footmisc}
\usepackage{color}
	\definecolor{maroon}{RGB}{128,0,0}
	\definecolor{navy}{RGB}{0,0,165}
	\definecolor{lb}{RGB}{150,165,190}
	\definecolor{MyWhite}{RGB}{255, 255, 255}
\usepackage{colortbl}

\usepackage{caption}
\usepackage{subcaption}
\usepackage{placeins}
\usepackage{nicefrac}
%\usepackage{chngcntr}

\newcommand{\bit}{\begin{itemize}}
\newcommand{\eit}{\end{itemize}}
\usepackage{listings}
\def\cN{\mathcal{N}}
\newtheorem{theorem}{Theorem}
\newtheorem{remark}[theorem]{Remark}
\usepackage{mdframed}
\newmdtheoremenv{assumption}{Assumption}

\usepackage{enumitem}


\begin{document}


%\onehalfspacing


\title{Monte Carlo Details for checking our simulation method}
%Beggar or Prosper-thy-Neighbour? \\
%Fiscal Spillovers Across Firms and Countries
\maketitle

\paragraph{Purpose:} The idea is to create data consistent with our estimation model that we can then run our estimation method on to see what we get.

\section{Shocks}

Let there be $J$ sectors. For each sector $j$ I want you to simulate two sets of shocks: a supply shock ($\nu_{j,t}$) and a demand shock ($\xi_{j,t}$). They should be simulated as follows:

\begin{align*}
	\nu_{j,t} &= \nu_{j,t-1} + v_{j,t}\\
	\text{with: }\;\; v_{j,t}&=\rho_j v_{j,t-1}+e_{j,t};\;\;e_{j,t}\sim N(0,\sigma_{e_j}^2)\\
	\xi_{j,t}&= \xi_{j,t-1} + \epsilon_{j,t}\\
	\text{with: }\;\;\epsilon_{j,t}&=p_j \epsilon_{j,t-1}+u_{j,t};\;\;u_{j,t}\sim N(0,\sigma_{u_j}^2)
\end{align*}

So both shocks are to be random walks that are themselves generated by persistent shocks. Parameters to be chosen by the user are $\{\rho_j,p_j\}_{j=1}^J$. To start with we can make them all the same.

One we have these shocks,let's put them into our model.

\section{Demand side of the model}

To start with let us simulate final demand. I motivated this project with final demand being Cobb-Douglas so let's stick with that. That means that expenditure shares are all constant and so the demand side of the model actually is fairly easy.

$$
	FD_{j,t} = \xi_{j,t}
$$

That is we assume that final demand is 100\% given by the demand shocks and 0\% due to supply shocks (or price changes). This is a special case we should think about relaxing later on but for now should make the simulations fairly easy.

\section{I-O data}

Let's use the data we get from the BEA. This should give us information on the share of intermediates used in production ($\Omega_{ij}$) for each sector as well as the share of all intermediates as a share of sales ($\gamma_j$).

\section{Supply Block}

Using the I-O we can map from final demand to sales fairly easily. This is done by:

$$
	\underline{Sales}_t = (I-\Omega' diag(\underline{\gamma}))^{-1}\underline{FD}_{t}
$$

where $\Omega$ is a matrix where the $(i,j)^{th}$ entry denotes the share of intermediates of type $j$ used in production by sector $i$ and $\underline{\gamma}$ is the share of all intermediate costs as a fraction of sales.

Next let's assume that production is given by:

$$
	Y_{j,t} = VA_{j,t}^{1-\gamma_j}M_{j,t}^{\gamma_{j}}
$$

where $VA_j$ is (real) value added and $M_j$ is intermediate quantities used in production.

Both have prices $P_{VA,j,t}$ and $P_{j,t}^{int}$ where if known we can use to calculate output prices:

$$
P_{j,t} = P_{VA,j,t}^{1-\gamma_j}(P_{j,t}^{int})^{\gamma_j}
$$

The issue is that I want to have an upward sloping supply curve so specifying $P_{VA,j,t}$ might be tricky.

So I want to specify an upward sloping supply curve as follows:

$$
	P_{VA,j,t} = Y_{j,t}^{\alpha_j}\nu_{j,t} 
$$

This means that as demand rises, so do costs and therefore prices. I am thinking we might want to set $\alpha_j>1$ but I don't care too much. To start with let's make it constant for all sectors.

Here is how we need to generate the fake data:

\begin{enumerate}
	\item Generate a guess of the vector of sectoral prices $\underline{p}_{t}^{(0)}$. To start with why not assume they are all 1.
	\item Given a guess of a vector of output prices $\underline{P}_t^{(n)}$, find the price of intermediates using $P_{j,t}^{(n)} = exp^{\sum_i \Omega_{ji}P_{i,t}^{(n)}}$
	\item Using the guess of output prices to find real output using the expenditure data: $\underline{Y}_t^{(n)} = \frac{\underline{Sales}_t}{\underline{P}_t^{(n)}}$ where the fraction represents \emph{element} division.
	\item Using real output, we can find the price of value added $P_{VA,j,t}^{(n)}$.
	\item Calculate a new output price $P_{j,t}^{(n+1)} = (P_{VA,j,t}^{(n)})^{1-\gamma_j}(P_{j,t}^{int,(n)})^{\gamma_j}$.
	\item Continue with steps (2) to (5) until the (largest absolute) change in prices is less than some tolerance that is small. I suggest $10^{-7}$. If this does not work well, we can always try Newton-Rhapson which should work fairly quickly. 
	\item Confirm that the equations of the model hold to a reasonable degree of accuracy.
\end{enumerate}

One that is done we have a time series of output prices $\underline{P}_t$ and final demands $\underline{FD}_{t}$ to use as data for our procedure. Then we can see how close we are getting to the (Monte-Carlo) truth.


\end{document}